% Academic CV (2 pages). Build: pdflatex Jose_Carter_CV_Academic_2026.tex
\documentclass[letterpaper,10pt]{article}
% Shared style for José Carter's résumés (matches the 2026 WDAS résumé look).
\usepackage[utf8]{inputenc}
\usepackage[T1]{fontenc}
\usepackage{lmodern}
\usepackage[english]{babel}
\usepackage[letterpaper,top=0.45in,bottom=0.45in,left=0.5in,right=0.5in]{geometry}
\usepackage{titlesec}
\usepackage[usenames,dvipsnames]{xcolor}
\usepackage{enumitem}
\usepackage{tabularx}
\usepackage[hidelinks,pdfusetitle]{hyperref}
\usepackage{microtype}

\definecolor{cvblue}{HTML}{1F3F73}
\pagestyle{empty}
\setlength{\parindent}{0pt}
\setlength{\tabcolsep}{0in}
\raggedright
\urlstyle{same}

\titleformat{\section}{\color{cvblue}\Large\bfseries}{}{0em}{}[\vspace{-2pt}\color{cvblue}\titlerule]
\titlespacing*{\section}{0pt}{4pt}{3pt}

\newcommand{\cvheader}[2]{%
  \begin{center}
    {\color{cvblue}\fontsize{22}{24}\selectfont\bfseries José Carter Arriagada}\\[2pt]
    {\small\bfseries #1}\\[2pt]
    {\small Santiago, Chile \textbar{} Chile: +56 9 8655 2408 \textbar{} U.S.: +1 276 221 3064}\\
    {\small \href{mailto:jose.carterx@gmail.com}{jose.carterx@gmail.com} \textbar{} \href{https://josecarter.dev}{josecarter.dev} #2}\\
    {\small \href{https://linkedin.com/in/jose-carter-arriagada}{linkedin.com/in/jose-carter-arriagada} \textbar{} \href{https://github.com/Cartterr}{github.com/Cartterr}}
  \end{center}
  \vspace{-10pt}
}

% \role{Title}{Dates}{Organization}{Location}
\newcommand{\role}[4]{%
  \vspace{1pt}
  \begin{tabular*}{\textwidth}{l@{\extracolsep{\fill}}r}
    \textbf{#1} & \small #2 \\
    \textit{\small #3} & \textit{\small #4} \\
  \end{tabular*}\vspace{-1pt}
}
\newenvironment{points}{\begin{itemize}[leftmargin=0.18in,topsep=0pt,itemsep=0pt,parsep=0pt,partopsep=0pt,label={\small\textbullet}]\small}{\end{itemize}\vspace{1pt}}
\newcommand{\entry}[2]{{\small\textbf{#1} --- #2}\par\vspace{1pt}}

\hypersetup{pdftitle={José Carter Arriagada - Academic CV 2026}}

\begin{document}
\cvheader{Research Software Engineer \textbar{} GPU Computing, Real-Time XR \& HPC}{}

\section{Research Interests}
{\small GPU and parallel computing for real-time and scientific workloads; real-time rendering and XR systems, including Gaussian splatting and neural rendering on consumer GPUs; performance engineering for simulation and physical AI; high-performance computing infrastructure for research.\par}

\section{Education}
\role{Pontificia Universidad Católica de Chile}{Mar 2020 -- Jul 2025}{B.S. in Computer Science Engineering; Major in Software Engineering; Minor in Data Science}{Santiago, Chile}
\begin{points}
  \item Graduated with Distinction. Coursework and teaching in Operating Systems, High-Performance Computing, Networks, and Extended Reality.
\end{points}

\section{Research Experience}
\role{Multiplatform Full Stack Software Engineer, XR Lab}{Sep 2026 -- Present}{DICTUC - Pontificia Universidad Católica de Chile, with Prof. Leonel Merino (full time)}{Santiago, Chile}
\begin{points}
  \item Lead software for \textbf{ANID project 13250116}, \emph{Enhancing Collaboration in Extended Reality}: a three-headset mixed-reality study tool for research on emotion and collaboration.
  \item Moved the study app to a native \textbf{VIVE Focus Vision} build (\textbf{89.5 fps}, passthrough, hand and eye tracking) and built a web panel that runs every headset from a PC or phone.
  \item Synchronized gaze, motion, and task events with \textbf{MindWare BioLab} physiological recordings to within a few milliseconds, with no turns lost on Wi-Fi drops, so sessions are usable for analysis.
  \item Built \textbf{Waves AR}, an iPhone and iPad app that teaches ocean waves on the coast (Swift, Metal FFT ocean, ARKit/RealityKit), and a wireless dev loop that cut an AR iteration from about 8 minutes to 4--7 seconds.
\end{points}

\role{Research Software Engineer}{Jul 2023 -- Jul 2025}{Politiktok (Fondecyt) - Pontificia Universidad Católica de Chile}{Santiago, Chile}
\begin{points}
  \item Led the engineering of a GPU-backed \textbf{Python/PostgreSQL} research platform processing \textbf{100,000+ records}, with about \textbf{10x faster} pipelines through parallelization and query optimization.
  \item Built reproducible ingestion, cleaning, transcription, feature-engineering, NLP, and visualization workflows for social-media research, supporting a peer-reviewed Springer publication.
\end{points}

\role{Simulation Engineer}{Jan 2023 -- Jul 2024}{Pontificia Universidad Católica de Chile - Geoscience}{Santiago, Chile}
\begin{points}
  \item Developed GPU-accelerated scientific software for tectonic-plate modeling and seismic-risk analysis, optimizing \textbf{CUDA} parallel algorithms for a \textbf{15x speedup}.
  \item Ran large simulations on Linux workstations and university servers, and processed and visualized the resulting data; presented the work at the SSA 2023 Annual Meeting.
\end{points}

\role{AI Systems Developer, UC Bridge Program}{Jan 2024 -- Mar 2024}{Drone Response - University of Notre Dame}{Notre Dame, IN, USA}
\begin{points}
  \item Selected among 200+ candidates for an international research program on NASA- and NSF-funded projects.
  \item Implemented mission-planning logic for autonomous multi-drone coordination in emergency response: Python route optimization (Hamiltonian and clustering approaches), MQTT event-driven communication, OpenAI API, and Java Spring services.
\end{points}

\section{Publications \& Presentations}
\entry{Journal article (2026)}{Co-author. \emph{Political Subjectification on TikTok: A Postdigital Mode of Inquiry for Tracing Affects Among Chilean Youth}. Postdigital Science and Education, Springer. \href{https://doi.org/10.1007/s42438-026-00638-4}{doi:10.1007/s42438-026-00638-4}. Built the technical pipeline for large-scale data mining, transcription, and analysis.}
\entry{Conference presentation (2023)}{Co-author. \emph{Interactions Between Megathrust and Adjacent Crustal Faults}. Seismological Society of America Annual Meeting, San Juan, Puerto Rico.}

\section{Teaching}
\role{Senior Teaching Assistant \& Technical Mentor}{Mar 2023 -- Jul 2025}{Department of Computer Science, Pontificia Universidad Católica de Chile}{Santiago, Chile}
\begin{points}
  \item 14+ teaching-assistant roles for 100+ students in \textbf{Operating Systems}, \textbf{High-Performance Computing}, Networks, Extended Reality, introductory computing, and software testing; mentored systems programming and optimization.
  \item Named Senior Teaching Assistant through the department's teaching-assistant career program (2025).
\end{points}
\role{Programming Instructor}{Jul 2024 -- Jul 2026}{Escuela Militar de Chile}{Santiago, Chile}
\begin{points}
  \item Designed and taught Python and AI-application courses, assessments, and materials; recognized in 2026 for dedication and professionalism.
\end{points}

\section{Industry \& Applied Experience}
\role{Freelance Software Engineer, Physical AI \& Robotics}{Sep 2026 -- Present}{TEK Chile (part time)}{Santiago, Chile}
\begin{points}
  \item Perception, control, and simulation software for a Dobot CR5S roll-packing cell and a Dobot Rover X1 inspection robot: Isaac Sim digital twin, YOLO26 vision, and Modbus/TCP control.
  \item Rover: 18 real patrol runs with no contact, video raised from 5 to about 20 fps, Pi latency cut from 125 to about 25 ms, and a 0.48 s emergency stop.
\end{points}
\role{Full Stack Developer}{Sep 2025 -- Sep 2026}{Dily - Automotive Fintech Platform}{Santiago, Chile}
\begin{points}
  \item End-to-end lending software with TypeScript/Node.js, Angular/Nx, PostgreSQL/MySQL, and AWS, using DDD and hexagonal architecture.
\end{points}
\role{Lead Engineer, Alerting Platform Migration}{2026}{GridWorks}{Santiago, Chile}
\begin{points}
  \item Led the staged rebuild of an industrial alerting platform (Next.js, Node.js, MQTT, PostgreSQL) without interrupting live workflows.
\end{points}
\role{Software Engineer}{Dec 2024 -- Jul 2025}{Flair - Commercial-Building Energy Optimization}{Santiago, Chile}
\begin{points}
  \item Python microservices and InfluxDB/DynamoDB time-series workflows on AWS for a platform tied to HVAC energy reductions of up to 50\%.
\end{points}

\section{Conferences, Programs \& Awards}
\entry{SC26 Early Career Program}{Selected participant with travel grant covering registration, airfare, and lodging; also matched in the SC26 Mentor-Protégé Program (Chicago, Nov 2026).}
\entry{ACM SIGGRAPH 2026 Student Volunteer}{Supported conference operations in Los Angeles (Jul 2026); visited Treyarch through the volunteer program.}
\entry{10th South American Space Generation Workshop}{Space Generation Advisory Council delegate (2026).}
\entry{KHIPU 2025}{Latin American Meeting in Artificial Intelligence.}

\section{Skills \& Languages}
{\small
\textbf{Programming:} Python, C++, CUDA, C\#, Swift, TypeScript/JavaScript, Java, SQL\\
\textbf{GPU \& Real-Time:} CUDA, Metal, Unity, OpenXR, ARKit/RealityKit, Isaac Sim, eye and hand tracking, profiling\\
\textbf{Data \& ML:} PyTorch, NLP, PostgreSQL, InfluxDB, data pipelines, scientific visualization\\
\textbf{Systems:} Linux, Docker, AWS, MQTT, Git, CI/CD\\
\textbf{Languages:} Spanish (native), English (advanced, TOEIC 905)\par}

\end{document}
